\chapter{Price, Value, and the Many Definitions of the Price}
\label{ch:price_value}
\chaptermark{Price and value}

Ask a newcomer what the price of a stock is, and the answer is almost
always ``the last trade price.'' That answer is not wrong, only
insufficient. On any real market, at any instant, several equally
valid numbers could be called ``the price,'' and a strategy that
picks the wrong one bleeds edge invisibly until the end-of-day
report. A market maker computing skew off the raw mid when one side
has a one-lot stub at top of book is quoting against noise, not fair
value. An execution algorithm benchmarking against a last trade from
four minutes ago looks good on paper and loses money in practice.
There is no single number; the skill this chapter builds is knowing
\emph{which number to use for which purpose}.

This chapter replaces the unqualified phrase ``the price'' with a
menu of well-defined price \emph{measures}, each with a formal
definition, a use case, and a failure mode: best bid, best ask,
spread, mid, microprice, volume-weighted average price, last-trade
price, and one Prosperity-specific construction called the Wall Mid.
By the end you will be able to compute every one on a given book and
state which answers the question you are actually asking.

\begin{backgroundbox}[Prerequisites]
From \cref{ch:lob}:
\begin{itemize}
  \item The \textbf{limit order book (LOB)} as two sorted lists of
  resting limit orders, and the definitions of \textbf{bid},
  \textbf{ask}, \textbf{spread}, and \textbf{mid} introduced there as
  quick working definitions. This chapter is the promised formal
  treatment.
  \item \textbf{Maker} vs \textbf{taker}: in every trade the maker
  was the resting side, the taker was the crossing side.
  \item \textbf{Tick size} and \textbf{lot size} as the price and
  quantity grids on which orders live.
\end{itemize}
From \cref{ch:participants}:
\begin{itemize}
  \item \textbf{Informed} vs \textbf{uninformed} traders, and the
  adverse-selection intuition: posted quotes are systematically hit
  by the side that knows more. This motivates why the microprice is
  typically a better short-horizon predictor than the mid.
\end{itemize}
No prior finance knowledge assumed beyond these.
\end{backgroundbox}

\section{Bid, ask, spread, mid: formal definitions}
\label{sec:bidaskmid_formal}

\Cref{ch:lob} introduced $\bidp$, $\askp$, $\spreads$, and $\midp$ as
working definitions in a single paragraph. The rest of this chapter
hangs off these four symbols, so we restate them formally here.

\begin{definition}[Best bid]
\index{bid price|textbf}
The \emph{best bid} $\bidp$ at time $t$ is the highest price among
all resting limit buy orders currently in the LOB:
\[
  \bidp \;\defeq\; \max \bigl\{\, \mathrm{price}(o) : o \in B_t \,\bigr\},
\]
where $B_t$ is the bid side of the book at time $t$. If $B_t$ is
empty, $\bidp$ is undefined.
\end{definition}

\begin{definition}[Best ask]
\index{ask price|textbf}
The \emph{best ask} (or \emph{best offer}) $\askp$ at time $t$ is the
lowest price among all resting limit sell orders currently in the
LOB:
\[
  \askp \;\defeq\; \min \bigl\{\, \mathrm{price}(o) : o \in A_t \,\bigr\}.
\]
If $A_t$ is empty, $\askp$ is undefined.
\end{definition}

\begin{definition}[Spread]
\index{bid-ask spread|textbf}\index{spread!bid-ask|textbf}
The \emph{bid--ask spread} $\spreads$ is the gap between the best ask
and the best bid, $\spreads \defeq \askp - \bidp$. On an uncrossed
book $\spreads \geq 0$; we have $\spreads = 0$ only in the degenerate
instant when the book is \emph{locked} mid-match, that is, when the
best bid and best ask are momentarily equal because the matching
engine has not yet cleared the crossing orders.
\end{definition}

\begin{definition}[Mid-price]
\index{mid price|textbf}
The \emph{mid-price} $\midp$ is the arithmetic average of the best
bid and the best ask:
\[
  \midp \;\defeq\; \tfrac{1}{2}\bigl(\bidp + \askp\bigr).
\]
\end{definition}

The mid is the most common ``summary'' of where a market is trading,
and $\midp$ appears in almost every subsequent chapter. It is not
``the price'' in any rigorous sense; it is a convenient average that
pretends the two sides of the book are symmetric. When they are not,
which is almost always, there are better choices.

\begin{pitfallbox}[``The mid is the price'' is wrong]
The mid is easy to compute, plot, and defend in a meeting. It also
implicitly asserts that the best bid and best ask carry \emph{equal}
information about where the next trade happens. This is almost never
true. If the bid is backed by $131$ lots and the ask by only $48$,
the ask clears first under a modest market buy, and the post-trade
mid sits strictly above the pre-trade mid, meaning the expected
next traded price was above the mid all along. The microprice
(\cref{sec:beyond_mid}) is the correction. Use the mid as an
\emph{input} to strategies, not as a belief about fair value.
\end{pitfallbox}

\section{Beyond the mid: three more price measures}
\label{sec:beyond_mid}

The mid treats the two sides as interchangeable. In reality, the
side with less resting quantity is more likely to be consumed next,
so the short-horizon expected price sits closer to the
\emph{opposite} side. The \textbf{microprice} makes this precise.
The \textbf{volume-weighted average price (VWAP)} serves a different
purpose: not ``what is fair value now'' but ``what was the average
traded price over some window.'' The \textbf{last-trade price} is
the cheapest summary, and the one most easily manipulated and most
prone to staleness. We take them in turn.

\subsection{The microprice}
\label{subsec:microprice}

\begin{definition}[Microprice]
Let $\bidp$ and $\askp$ be the best bid and best ask, and let $Q_b$
and $Q_a$ be the total resting quantities at those two top-of-book
levels. The \emph{microprice} is
\[
  \mprice \;\defeq\; \frac{Q_a \cdot \bidp + Q_b \cdot \askp}{Q_a + Q_b}.
\]
It is a weighted average of the two top quotes in which each side's
\emph{own} quantity is the \emph{opposite} side's weight.
\end{definition}

The cross-weighting is the whole point. The side with smaller size
gets the larger weight: if $Q_a$ is tiny, the microprice leans hard
toward $\askp$, because an incoming market buy of modest size
consumes $Q_a$ and lifts the ask. Symmetrically, small $Q_b$ pulls
it toward $\bidp$. When $Q_a = Q_b$ the weights are equal and
$\mprice = \midp$; in every other case $\mprice$ and $\midp$
disagree, and the disagreement is informative.

Bouchaud, Bonart, Donier, and Gould \citeBouchaud{} give the formal
treatment and show the microprice is a markedly better short-horizon
predictor of the next mid movement than the mid itself: it
anticipates the mid rather than lagging it. The underlying reason is order-flow
imbalance, developed in \cref{ch:order_flow}.

\begin{intuitionbox}
The side with less size breaks first: a smaller buffer against the
next market order means a higher probability that side gets wiped
and its top quote moves. If you believe that (and the data in
\cref{ch:order_flow} say you should), then the near-future price
sits closer to the \emph{other} side. The microprice is the cheapest
implementation of that belief: one weighted average, no model, no
parameter to fit. Whenever a market-making chapter refers to ``fair
value'' as an input to a quoting formula, assume the author means
microprice unless stated otherwise.
\end{intuitionbox}

\subsection{VWAP}
\label{subsec:vwap}

\begin{definition}[Volume-weighted average price (VWAP)]
Given a time window $[t_0, t_1]$ over which $n$ trades
$(p_1, q_1), \ldots, (p_n, q_n)$ have printed on a product, the
\emph{volume-weighted average price} over that window is
\[
  \mathrm{VWAP}_{[t_0,t_1]} \;\defeq\; \frac{\sum_{i=1}^n p_i \, q_i}{\sum_{i=1}^n q_i}.
\]
\end{definition}

VWAP is not a fair-value estimate; it is a \emph{benchmark}. The
standard use case is execution measurement. Suppose a pension fund
buys a million shares over a trading day. The fund does not care
about the current mid; it cares whether its average acquisition
price beat or missed the day's VWAP. Beating VWAP by a basis point
on a million-share parent order is a valuable sell-side
relationship; missing by five basis points is an excuse to find a
different broker. (One \emph{basis point}, or \emph{bp}, is one
hundredth of one percent: $1\,\mathrm{bp} = 0.01\% = 0.0001$. It is
the standard unit for quoting small return differences on a trading
floor.) Almgren and Chriss's optimal execution framework
(\cref{ch:almgren_chriss}) is ultimately an argument about trading
off market impact against the risk of drifting from a VWAP-like
benchmark.

Two practical notes. First, VWAP depends on the window: ``day
VWAP'' is the common version, but ``trailing one-hour VWAP'' and
``VWAP since open'' appear routinely on execution dashboards.
Second, VWAP is backward-looking (it exists only after the trades
print), so it cannot be a forward-looking fair value. Do not
confuse the two.

\subsection{Last-trade price}
\label{subsec:last_trade}

\begin{definition}[Last-trade price]
\index{last traded price|textbf}
The \emph{last-trade price} $p^{\mathrm{last}}_t$ is the price of
the most recent trade to have printed on the product strictly before
time $t$. If no trades have yet occurred in the current session,
$p^{\mathrm{last}}_t$ is undefined.
\end{definition}

The last-trade price is the folk-meaning of ``the price'': what
retail websites display as ``current'' and what the closing-auction
print becomes for end-of-day accounting. Its use case is narrow but
important: \textbf{marking a position to market}. At session end
every firm must assign a dollar value to every position it holds
(for profit-and-loss reporting, margin, and risk limits), and that
dollar value has to come from \emph{somewhere}. For liquid products trading every
second, the last-trade price is cheap, well-defined, and consistent
across firms. For illiquid products trading every half hour, it is
stale and can be far from the current market; the mid, or an agreed
closing-auction print, is preferred.

Its weakness: a single number from a noisy stream. One small trade
at an odd price moves it; one unwatched cancelled trade leaves it
stale for an hour. It is also the cheapest price to manipulate near
the close (the regulatory phrase \emph{``marking the close''}
refers to exactly this), which is why closing auctions exist as a
separate mechanism.

\section{Wall Mid: a Prosperity-specific fair value}
\label{sec:wall_mid}

Prosperity publishes a limited-depth snapshot each tick, typically
three or four levels per side. Most visible flow comes from a few
bot market makers whose quotes sit on the \emph{outside} of the
visible book; the levels \emph{between} top of book and those outer
bot quotes are transient orders from competing human-written
algorithms, including one-lot ``stub'' orders posted for signalling
or to momentarily narrow the top. Using raw top-of-book as a
fair-value proxy therefore reacts to every skirmish there. The
Frankfurt Hedgehogs noticed that the \emph{outer} visible levels
(the lowest bid and highest ask in each snapshot) were
consistently posted by the same designated market-maker bots,
quoting rounded versions of the true underlying price (typically
$\pm 2$ ticks). They named those outer levels the \emph{bid wall}
and \emph{ask wall}, and defined the \textbf{Wall Mid} as their
midpoint.

\begin{definition}[Bid wall, ask wall, Wall Mid]
Given the visible depth of a Prosperity LOB snapshot, the \emph{bid
wall} is the bid price of the deepest (lowest-priced) visible bid
level, and the \emph{ask wall} is the ask price of the deepest
(highest-priced) visible ask level:
\begin{align*}
  p^{\mathrm{bid,wall}} &\;\defeq\; \min \{\, p : p \in \text{visible bid levels}\, \},\\
  p^{\mathrm{ask,wall}} &\;\defeq\; \max \{\, p : p \in \text{visible ask levels}\, \}.
\end{align*}
The \emph{Wall Mid} is the midpoint between the two walls:
\[
  \wallmid \;\defeq\; \tfrac{1}{2}\bigl(p^{\mathrm{bid,wall}} + p^{\mathrm{ask,wall}}\bigr).
\]
\end{definition}

The construction is deliberately crude: a one-line, parameter-free
formula computable in constant time from any snapshot. Its only
claim is that on the Prosperity platform the outer visible levels
were posted by informed market-maker bots whose quotes were
near-unbiased estimates of a hidden ``true'' price, so averaging the
two outer levels gave a far more stable estimate than any
top-of-book quantity.

\begin{prosperitybox}
\textbf{Wall Mid: the Frankfurt Hedgehogs' signature construction.}
Below is the exact helper used by the 2nd-place Frankfurt Hedgehogs
algorithm in Prosperity~3 to compute the Wall Mid from a
\texttt{TradingState}. \texttt{mkt\_buy\_orders} and
\texttt{mkt\_sell\_orders} are price$\to$size dictionaries for the
two sides of the visible depth; \texttt{bid\_wall} is the
\emph{lowest} visible bid and \texttt{ask\_wall} the \emph{highest}
visible ask, matching the definition above. The defensive
\texttt{try/except} blocks handle ticks where one side is empty.

\begin{lstlisting}[language=Python, caption={Wall Mid helper from \texttt{FrankfurtHedgehogs\_polished.py}, lines 153--166, verbatim.}]
def get_walls(self):

    bid_wall = wall_mid = ask_wall = None

    try: bid_wall = min([x for x,_ in self.mkt_buy_orders.items()])
    except: pass

    try: ask_wall = max([x for x,_ in self.mkt_sell_orders.items()])
    except: pass

    try: wall_mid = (bid_wall + ask_wall) / 2
    except: pass

    return bid_wall, wall_mid, ask_wall
\end{lstlisting}

Every strategy in the Frankfurt Hedgehogs writeup
(\texttt{imc-prosperity-3/README.md}) feeds its taking decisions and
quote skew off this single \texttt{wall\_mid} number. It is the one
custom construction separating their 2nd-place algorithm from a
generic top-of-book market maker, and the single piece of
Prosperity-specific insight most worth stealing for a 2026 entry.
\end{prosperitybox}

\section{A running numerical example}
\label{sec:running_example}

Take the same four-level toy LOB carried through \cref{ch:lob},
reproduced with sizes made explicit, since for this chapter every
size matters.

\begin{center}
\begin{tabular}{cc|cc}
\toprule
\multicolumn{2}{c|}{\textbf{BID SIDE}} & \multicolumn{2}{c}{\textbf{ASK SIDE}} \\
Qty & Price & Price & Qty \\
\midrule
$131$ & $102.54$ & $103.23$ & $48$  \\
$32$  & $101.87$ & $103.98$ & $84$  \\
$293$ & $101.48$ & $104.17$ & $38$  \\
$65$  & $101.10$ & $104.75$ & $127$ \\
\bottomrule
\end{tabular}
\end{center}

Assume ten trades have just printed in the VWAP window:

\begin{center}
\begin{tabular}{lcc}
\toprule
Trade & Price & Qty \\
\midrule
$T_1, T_2, T_3, T_4, T_5$       & $102.80$ & $10$ each \\
$T_6, T_7, T_8$                 & $103.00$ & $20$ each \\
$T_9, T_{10}$                   & $103.20$ & $15$ each \\
\bottomrule
\end{tabular}
\end{center}

The most recent trade $T_{10}$ printed at $103.20$. We now compute
every price measure on this one book.

\paragraph{Bid, ask, spread, mid.} By inspection, $\bidp = 102.54$
and $\askp = 103.23$, so
\begin{align*}
  \spreads &\;=\; \askp - \bidp \;=\; 103.23 - 102.54 \;=\; 0.69,\\
  \midp    &\;=\; \tfrac{1}{2}(102.54 + 103.23) \;=\; 102.885.
\end{align*}

\paragraph{Microprice.} Top-of-book sizes are $Q_b = 131$, $Q_a = 48$,
so $\askp$ gets weight $Q_b/(Q_a+Q_b)$ and $\bidp$ weight $Q_a/(Q_a+Q_b)$:
\[
  \mprice
  \;=\; \frac{Q_a \cdot \bidp + Q_b \cdot \askp}{Q_a + Q_b}
  \;=\; \frac{48 \cdot 102.54 + 131 \cdot 103.23}{48 + 131}.
\]
Numerator $48 \cdot 102.54 + 131 \cdot 103.23 = 4{,}921.92 + 13{,}523.13 = 18{,}445.05$;
denominator $179$; so
\[
  \mprice \;=\; \frac{18{,}445.05}{179} \;\approx\; 103.0450.
\]
The microprice sits $103.0450 - 102.8850 = 0.1600$ \emph{above} the
mid. That is the signal: because the bid is stacked roughly
$131/48 \approx 2.7$ times deeper than the ask, the next trade is
more likely to be a buy lifting the ask, and the microprice expects
the top of book to drift upward.

\paragraph{VWAP.} Using the ten trades in the window, the
size-weighted numerator is
\begin{align*}
  \sum_i p_i q_i
  &\;=\; \underbrace{50 \cdot 102.80}_{T_1..T_5}
       + \underbrace{60 \cdot 103.00}_{T_6..T_8}
       + \underbrace{30 \cdot 103.20}_{T_9,T_{10}} \\
  &\;=\; 5{,}140 + 6{,}180 + 3{,}096
  \;=\; 14{,}416,
\end{align*}
and the total traded quantity is $\sum_i q_i = 50 + 60 + 30 = 140$,
giving
\[
  \mathrm{VWAP} \;=\; \frac{14{,}416}{140} \;\approx\; 102.9714.
\]

\paragraph{Last-trade price.} $T_{10}$ printed at $103.20$, so
$p^{\mathrm{last}} = 103.20$.

\paragraph{Wall Mid.} The bid wall is the lowest visible bid,
$p^{\mathrm{bid,wall}} = 101.10$; the ask wall the highest visible
ask, $p^{\mathrm{ask,wall}} = 104.75$. So
\[
  \wallmid \;=\; \tfrac{1}{2}(101.10 + 104.75) \;=\; \tfrac{1}{2} \cdot 205.85 \;=\; 102.925.
\]

\paragraph{Summary table.} All six price measures on one book:

\begin{center}
\begin{tabular}{llr}
\toprule
Symbol & Measure & Value \\
\midrule
$\bidp$                 & Best bid                          & $102.54$   \\
$\askp$                 & Best ask                          & $103.23$   \\
$\spreads$              & Spread                            & $0.69$     \\
$\midp$                 & Mid-price                         & $102.885$  \\
$\mprice$               & Microprice $(Q_b{=}131, Q_a{=}48)$ & $103.0450$ \\
$\mathrm{VWAP}$         & VWAP (10-trade window)            & $102.9714$ \\
$p^{\mathrm{last}}$     & Last-trade price                  & $103.20$   \\
$\wallmid$              & Wall Mid                          & $102.925$  \\
\bottomrule
\end{tabular}
\end{center}

Six numbers, all nominally near ``one hundred and three,'' spread
across a range of $103.20 - 102.54 = 0.66$, almost exactly the
quoted spread. This dispersion is the point. Each number answers a
different question, and confusing them is the cheapest way a new
quant loses money.

\begin{keyfactbox}[Which price when?]
\begin{itemize}
  \item \textbf{Quoting a passive limit order} (market making) $\to$
  $\mprice$ or $\wallmid$, not $\midp$. Use the one the book's
  top-of-book noise behaviour prefers.
  \item \textbf{Marking a position to market} at session end $\to$
  $p^{\mathrm{last}}$ for liquid products, $\midp$ for illiquid
  ones where the last trade is stale.
  \item \textbf{Benchmarking an execution} over a time window $\to$
  $\mathrm{VWAP}$ over that window. Nothing else.
  \item \textbf{Estimating ``fair value''} as an input to a taking or
  hedging decision $\to$ $\mprice$ on a real exchange,
  $\wallmid$ on Prosperity.
\end{itemize}
Never use a single one of the six numbers for all four purposes.
\end{keyfactbox}

\section{Fair value, informally}
\label{sec:fair_value}

Every price measure in this chapter is a \emph{statistic}: a
deterministic function of the LOB or the recent trade tape. None
estimates a latent quantity. The word quants use for that latent
quantity is \textbf{fair value}; since the term is thrown around
casually on trading desks, we state clearly what it means.

\begin{definition}[Fair value, informally]
\index{fundamental value|textbf}
The \emph{fair value} of an asset at time $t$ is, informally, the
price at which the asset would trade in a frictionless market in
which all currently-available information had already been fully
incorporated into prices. It is not an observable number; it is a
model-dependent \emph{estimate}, and different models will give
different fair values for the same asset at the same instant.
\end{definition}

The informal definition suffices here. Two rigorous instances appear
later. \Cref{ch:spreads_roll} covers Roll's 1984 bid--ask spread
model, which defines the \emph{efficient price} as a latent random
walk of unobservable innovations and extracts it from trade-price
autocovariance alone. \Cref{ch:black_scholes} covers Black--Scholes,
where the fair value of a European option is the unique
no-arbitrage price implied by the underlying's dynamics, rigorous
once inputs are fixed but inheriting all the fragility of those
inputs. In both cases fair value is the output of a model, and the
model is only as good as its inputs and assumptions.

\begin{pitfallbox}[The fair-value trap]
The most seductive mistake in quant trading is to build a
fair-value model, believe its output, and trade size against it. A
fair-value model can be wrong in at least three uncorrelated ways:
(i) its functional form is misspecified; (ii) its inputs are stale,
noisy, or wrong; (iii) its calibration data come from a regime that
no longer applies. Any one alone produces a confident estimate
dollars away from where the market trades. \Cref{ch:backtesting}
treats the calibration trap: a model fit on a window with
survivorship bias or look-ahead agrees with that history by
construction and disagrees with live markets exactly when it
matters. The discipline is to treat every fair-value number as one
signal among many and to size positions against the
\emph{dispersion} of candidate fair values, not the one you like
best.
\end{pitfallbox}

\section{Key facts to memorise}
\label{sec:ch3_keyfacts}

\begin{itemize}
  \item $\bidp$ and $\askp$ are the top of the two sides of the LOB;
  $\spreads = \askp - \bidp$ and $\midp = (\bidp + \askp)/2$. These
  four numbers are the minimum summary of a market's state, but
  \emph{none} of them is ``the price.''
  \item The \textbf{microprice} is the cross-size-weighted mid:
  $\mprice = (Q_a \bidp + Q_b \askp)/(Q_a + Q_b)$. The side with
  less size gets the larger weight, because that side will break
  first.
  \item The microprice is a short-horizon predictor of the next mid
  move; its motivating quantity is \emph{order-flow imbalance},
  developed in \cref{ch:order_flow}.
  \item \textbf{VWAP} is a \emph{backward-looking} benchmark over a
  window, not a fair-value estimate. It is the standard yardstick
  for institutional execution quality and the reference point for
  \cref{ch:almgren_chriss}.
  \item The \textbf{last-trade price} is the cheapest summary of
  where the market is trading, used for end-of-day marking in liquid
  products, but it is noisy, stale in illiquid products, and
  manipulable near the close.
  \item The \textbf{Wall Mid} $\wallmid$ is the Frankfurt Hedgehogs'
  Prosperity-specific fair-value proxy: the midpoint of the outermost
  visible bid and ask levels, which in Prosperity correspond to the
  designated-market-maker bot quotes.
  \item \textbf{Fair value} is not a statistic; it is a
  model-dependent estimate of a latent frictionless price, and
  different models give different numbers. Every model is only as
  good as its inputs and its calibration data
  (\cref{ch:backtesting}).
  \item The one-sentence rule: \emph{there is no single price, only
  price measures, and the skill is knowing which measure answers
  which question.}
\end{itemize}
